\documentclass[11pt]{article}

\usepackage[margin=1in]{geometry}
\usepackage{amsmath, amssymb, array, booktabs}
\usepackage[T1]{fontenc}
\usepackage{lmodern}
\usepackage{hyperref}

\title{Step-by-Step Derivatives of the Lagrange and Hermite Shape Functions}
\author{Companion notes for \texttt{FEMHermiteBeamRegion.py}}
\date{}

\newcommand{\dd}{\,d}
\newcommand{\pd}[2]{\frac{\partial #1}{\partial #2}}
\newcommand{\od}[2]{\frac{d #1}{d #2}}
\newcommand{\xip}{\xi}
\newcommand{\etap}{\eta}
\newcommand{\zetap}{\zeta}
\newcommand{\Hs}{H_{\mathrm{scale}}}

\begin{document}
\maketitle

\section*{1. The derivative rules used}

Every derivative in the code comes from a small set of Calculus 1 rules.

\begin{center}
\begin{tabular}{>{\raggedright\arraybackslash}p{0.24\linewidth}
                >{\raggedright\arraybackslash}p{0.33\linewidth}
                >{\raggedright\arraybackslash}p{0.34\linewidth}}
\toprule
\textbf{Rule} & \textbf{Formula} & \textbf{How it appears here} \\
\midrule
Constant multiple rule
& $\displaystyle \od{}{s}\left(c f(s)\right)=c f'(s)$
& The factors $\frac18$, $\frac14$, and $\Hs$ stay outside the derivative. \\
\addlinespace
Sum and difference rule
& $\displaystyle \od{}{s}\left(f(s)\pm g(s)\right)=f'(s)\pm g'(s)$
& The Hermite cubic polynomials are differentiated term by term. \\
\addlinespace
Power rule
& $\displaystyle \od{}{s}\left(s^n\right)=n s^{n-1}$
& For example, $\od{}{s}(s^3)=3s^2$ and $\od{}{s}(s^2)=2s$. \\
\addlinespace
Partial derivative rule
& When differentiating with respect to $\xi$, treat $\eta$ and $\zeta$ as constants.
& In Hex8 functions, only the factor containing the chosen variable changes. \\
\addlinespace
Product rule
& $\displaystyle \od{}{s}(fg)=f'g+fg'$
& Used for the tensor-product Hermite functions. Since most factors do not depend on the variable, many product-rule terms become zero. \\
\bottomrule
\end{tabular}
\end{center}

\section*{2. Lagrange Hex8 shape functions}

The eight trilinear Lagrange shape functions can be written compactly as
\[
N_i(\xi,\eta,\zeta)
=
\frac18
(1+a_i\xi)(1+b_i\eta)(1+c_i\zeta),
\]
where each node has signs $a_i,b_i,c_i\in\{-1,+1\}$.
For the node ordering used in the code,
\[
\begin{array}{c|ccc}
i & a_i & b_i & c_i\\
\hline
1 & -1 & -1 & -1\\
2 & +1 & -1 & -1\\
3 & +1 & +1 & -1\\
4 & -1 & +1 & -1\\
5 & -1 & -1 & +1\\
6 & +1 & -1 & +1\\
7 & +1 & +1 & +1\\
8 & -1 & +1 & +1
\end{array}
\]

\subsection*{2.1 General derivative pattern}

Take the partial derivative with respect to $\xi$:
\[
\pd{N_i}{\xi}
=
\pd{}{\xi}
\left[
\frac18(1+a_i\xi)(1+b_i\eta)(1+c_i\zeta)
\right].
\]
By the constant multiple rule,
\[
\pd{N_i}{\xi}
=
\frac18
\pd{}{\xi}
\left[
(1+a_i\xi)(1+b_i\eta)(1+c_i\zeta)
\right].
\]
While differentiating with respect to $\xi$, the factors
$(1+b_i\eta)$ and $(1+c_i\zeta)$ are constants. Therefore,
\[
\pd{N_i}{\xi}
=
\frac18
\left[
\pd{}{\xi}(1+a_i\xi)
\right]
(1+b_i\eta)(1+c_i\zeta).
\]
Since $\pd{}{\xi}(1+a_i\xi)=a_i$,
\[
\boxed{
\pd{N_i}{\xi}
=
\frac18 a_i(1+b_i\eta)(1+c_i\zeta)
}.
\]

The same logic gives
\[
\boxed{
\pd{N_i}{\eta}
=
\frac18 b_i(1+a_i\xi)(1+c_i\zeta)
},
\qquad
\boxed{
\pd{N_i}{\zeta}
=
\frac18 c_i(1+a_i\xi)(1+b_i\eta)
}.
\]

\subsection*{2.2 Example: node 1}

The first shape function is
\[
N_1=\frac18(1-\xi)(1-\eta)(1-\zeta).
\]

Differentiate with respect to $\xi$:
\[
\pd{N_1}{\xi}
=
\frac18
\left[
\pd{}{\xi}(1-\xi)
\right]
(1-\eta)(1-\zeta)
=
\frac18(-1)(1-\eta)(1-\zeta).
\]
So
\[
\boxed{\pd{N_1}{\xi}=-\frac18(1-\eta)(1-\zeta)}.
\]

Differentiate with respect to $\eta$:
\[
\pd{N_1}{\eta}
=
\frac18(1-\xi)
\left[
\pd{}{\eta}(1-\eta)
\right]
(1-\zeta)
=
\frac18(1-\xi)(-1)(1-\zeta),
\]
so
\[
\boxed{\pd{N_1}{\eta}=-\frac18(1-\xi)(1-\zeta)}.
\]

Differentiate with respect to $\zeta$:
\[
\pd{N_1}{\zeta}
=
\frac18(1-\xi)(1-\eta)
\left[
\pd{}{\zeta}(1-\zeta)
\right]
=
\frac18(1-\xi)(1-\eta)(-1),
\]
so
\[
\boxed{\pd{N_1}{\zeta}=-\frac18(1-\xi)(1-\eta)}.
\]

\subsection*{2.3 All Hex8 derivatives}

The code stores these derivatives in a matrix whose rows are
\[
\left[
\pd{}{\xi},\quad
\pd{}{\eta},\quad
\pd{}{\zeta}
\right].
\]
For the eight shape functions:
\[
\pd{\mathbf N}{\xi}
=
\frac18
\begin{bmatrix}
-(1-\eta)(1-\zeta)\\
 (1-\eta)(1-\zeta)\\
 (1+\eta)(1-\zeta)\\
-(1+\eta)(1-\zeta)\\
-(1-\eta)(1+\zeta)\\
 (1-\eta)(1+\zeta)\\
 (1+\eta)(1+\zeta)\\
-(1+\eta)(1+\zeta)
\end{bmatrix}^{T},
\]
\[
\pd{\mathbf N}{\eta}
=
\frac18
\begin{bmatrix}
-(1-\xi)(1-\zeta)\\
-(1+\xi)(1-\zeta)\\
 (1+\xi)(1-\zeta)\\
 (1-\xi)(1-\zeta)\\
-(1-\xi)(1+\zeta)\\
-(1+\xi)(1+\zeta)\\
 (1+\xi)(1+\zeta)\\
 (1-\xi)(1+\zeta)
\end{bmatrix}^{T},
\]
\[
\pd{\mathbf N}{\zeta}
=
\frac18
\begin{bmatrix}
-(1-\xi)(1-\eta)\\
-(1+\xi)(1-\eta)\\
-(1+\xi)(1+\eta)\\
-(1-\xi)(1+\eta)\\
 (1-\xi)(1-\eta)\\
 (1+\xi)(1-\eta)\\
 (1+\xi)(1+\eta)\\
 (1-\xi)(1+\eta)
\end{bmatrix}^{T}.
\]

\section*{3. One-dimensional Hermite shape functions}

The four one-dimensional Hermite basis functions in the code are
\[
h_1(s)=\frac14(2-3s+s^3),
\]
\[
h_2(s)=\frac14(1-s-s^2+s^3),
\]
\[
h_3(s)=\frac14(2+3s-s^3),
\]
\[
h_4(s)=\frac14(-1-s+s^2+s^3).
\]
The function \texttt{hermite\_H\_functions} returns
\[
\Hs
\begin{bmatrix}
h_1 & h_2 & h_3 & h_4
\end{bmatrix}^{T}.
\]
Since $\Hs$ is constant with respect to $s$, the derivative is simply
$\Hs$ times the derivative of each $h_i$.

\subsection*{3.1 Derivative of $h_1$}

\[
h_1(s)=\frac14(2-3s+s^3).
\]
Use the constant multiple rule:
\[
h_1'(s)=
\frac14\od{}{s}(2-3s+s^3).
\]
Differentiate term by term:
\[
\od{}{s}(2)=0,\qquad
\od{}{s}(-3s)=-3,\qquad
\od{}{s}(s^3)=3s^2.
\]
Therefore
\[
\boxed{
h_1'(s)=\frac14(-3+3s^2)
}.
\]

\subsection*{3.2 Derivative of $h_2$}

\[
h_2(s)=\frac14(1-s-s^2+s^3).
\]
Then
\[
h_2'(s)=
\frac14
\left[
0-1-2s+3s^2
\right],
\]
so
\[
\boxed{
h_2'(s)=\frac14(-1-2s+3s^2)
}.
\]

\subsection*{3.3 Derivative of $h_3$}

\[
h_3(s)=\frac14(2+3s-s^3).
\]
Then
\[
h_3'(s)=
\frac14
\left[
0+3-3s^2
\right],
\]
so
\[
\boxed{
h_3'(s)=\frac14(3-3s^2)
}.
\]

\subsection*{3.4 Derivative of $h_4$}

\[
h_4(s)=\frac14(-1-s+s^2+s^3).
\]
Then
\[
h_4'(s)=
\frac14
\left[
0-1+2s+3s^2
\right],
\]
so
\[
\boxed{
h_4'(s)=\frac14(-1+2s+3s^2)
}.
\]

Thus the code's \texttt{dhermite\_H\_functions} returns
\[
\boxed{
\Hs
\begin{bmatrix}
\frac14(-3+3s^2)\\
\frac14(-1-2s+3s^2)\\
\frac14(3-3s^2)\\
\frac14(-1+2s+3s^2)
\end{bmatrix}
}.
\]

\section*{4. How the 3D Hermite derivatives are assembled}

The method \texttt{\_node\_axis\_shape} chooses one value-like Hermite function
and one rotation-like Hermite function for each coordinate direction.
For a negative node sign it chooses
\[
v=h_1,\qquad r=h_2,
\]
and for a positive node sign it chooses
\[
v=h_3,\qquad r=h_4.
\]

For a node with signs $(s_x,s_y,s_z)$, the code evaluates
\[
v_x=v(\xi),\quad v_y=v(\eta),\quad v_z=v(\zeta),
\]
and their derivatives
\[
v_x'=\od{v_x}{\xi},\quad
v_y'=\od{v_y}{\eta},\quad
v_z'=\od{v_z}{\zeta}.
\]

\subsection*{4.1 Translation-like Hermite shape}

The scalar translation-like Hermite shape for one node is
\[
N_H=v_xv_yv_z.
\]
Differentiate with respect to $\xi$:
\[
\pd{N_H}{\xi}
=
\pd{}{\xi}(v_xv_yv_z).
\]
Only $v_x$ depends on $\xi$. The other two factors are constants during this
partial derivative, so the product rule reduces to
\[
\boxed{
\pd{N_H}{\xi}=v_x'v_yv_z
}.
\]
Likewise,
\[
\boxed{
\pd{N_H}{\eta}=v_xv_y'v_z
},
\qquad
\boxed{
\pd{N_H}{\zeta}=v_xv_yv_z'
}.
\]
These are the three entries assigned to \texttt{dNH[:, n]}:
\[
\text{\texttt{dNH[:, n]}}=
\begin{bmatrix}
v_x'v_yv_z\\
v_xv_y'v_z\\
v_xv_yv_z'
\end{bmatrix}.
\]

\subsection*{4.2 Rotation-like Hermite shapes}

The code creates three rotation-like shapes per node:
\[
R_x=r_xv_yv_z,\qquad
R_y=v_xr_yv_z,\qquad
R_z=v_xv_yr_z.
\]

For $R_x$, the partial derivatives are
\[
\boxed{
\pd{R_x}{\xi}=r_x'v_yv_z
},
\qquad
\boxed{
\pd{R_x}{\eta}=r_xv_y'v_z
},
\qquad
\boxed{
\pd{R_x}{\zeta}=r_xv_yv_z'
}.
\]

For $R_y$:
\[
\boxed{
\pd{R_y}{\xi}=v_x'r_yv_z
},
\qquad
\boxed{
\pd{R_y}{\eta}=v_xr_y'v_z
},
\qquad
\boxed{
\pd{R_y}{\zeta}=v_xr_yv_z'
}.
\]

For $R_z$:
\[
\boxed{
\pd{R_z}{\xi}=v_x'v_yr_z
},
\qquad
\boxed{
\pd{R_z}{\eta}=v_xv_y'r_z
},
\qquad
\boxed{
\pd{R_z}{\zeta}=v_xv_yr_z'
}.
\]

That is exactly what appears in the code:
\[
\text{\texttt{dRH[:, 0, n]}}=
\begin{bmatrix}
r_x'v_yv_z\\
r_xv_y'v_z\\
r_xv_yv_z'
\end{bmatrix},
\quad
\text{\texttt{dRH[:, 1, n]}}=
\begin{bmatrix}
v_x'r_yv_z\\
v_xr_y'v_z\\
v_xr_yv_z'
\end{bmatrix},
\quad
\text{\texttt{dRH[:, 2, n]}}=
\begin{bmatrix}
v_x'v_yr_z\\
v_xv_y'r_z\\
v_xv_yr_z'
\end{bmatrix}.
\]

\section*{5. Connection to the Python code}

The Lagrange derivative routine returns
\[
\text{\texttt{dN\_dn}}
=
\begin{bmatrix}
\pd{N_1}{\xi} & \cdots & \pd{N_8}{\xi}\\
\pd{N_1}{\eta} & \cdots & \pd{N_8}{\eta}\\
\pd{N_1}{\zeta} & \cdots & \pd{N_8}{\zeta}
\end{bmatrix}.
\]
In the source code, the array is first written node-by-node as rows of
\[
\left[
\pd{N_i}{\xi},\quad
\pd{N_i}{\eta},\quad
\pd{N_i}{\zeta}
\right],
\]
then transposed with \texttt{.T}. That transpose turns it into the
``derivative direction by node'' layout shown above.

The Hermite derivative routine follows the same mathematical idea, but the
shapes are built from one-dimensional cubic Hermite functions. First the code
differentiates the one-dimensional cubic polynomials. Then it uses the product
rule to differentiate the tensor products in $\xi$, $\eta$, and $\zeta$.

\end{document}
